% !TeX root = ../../Book3.tex
% 纲要标题：### G.3 非负性、平方和与 Hilbert 第17问题

\section{非负性、平方和与 Hilbert 第17问题}
\label{b3:sec:G03}

平方和恒等式可以直接证明非负性，但多元非负多项式未必是多项式的平方和。允许有理函数及其分母后，能够得到更一般的正性表示；分母的零点仍需另行检查。

% 纲要细目（与计划纲要同步）：
% * 区别非负多项式、多项式平方和与有理函数平方和；Motzkin 非平方和证明明确最高次平方不能相消
% * 通过有理函数域赋序、正锥扩张、实闭包和量词消去完整证明 Artin 定理，不借实零点定理
% * 有理函数证书先清分母成为多项式恒等式；系数域、分母零点及原分式的定义域分别说明
% #### G.3.1 非负多项式、多项式平方和与有理函数平方和
% #### G.3.2 Motzkin 多项式与 Artin 定理
% #### G.3.3 平方和表示的系数域、分母与定义域

\subsection{非负多项式、多项式平方和与有理函数平方和}
\label{b3:subsec:G03:01}

设 \(n\ge1\)，\(f\in\mathbb R[x_1,\ldots,x_n]\)。若对每个 \(a\in\mathbb R^n\) 都有 \(f(a)\ge0\)，则称 \(f\) 为非负多项式。若有恒等式 \(f=\sum_jq_j^2\)，其中各 \(q_j\in\mathbb R[x_1,\ldots,x_n]\)，则称 \(f\) 为\term{平方和}[sum of squares]。两者的差别，是逐点的符号条件能否由某种有限代数恒等式表现出来。

\ProseParagraphBreak
这三个条件之间首先有如下关系：
\[
\begin{aligned}
\text{多项式平方和}&\ \Longrightarrow\ \text{在}\;\mathbb R^n\;\text{上非负},\\
\text{在}\;\mathbb R^n\;\text{上非负}&\ \xLongrightarrow{\text{Artin 定理}}\ \text{有理函数平方和}.
\end{aligned}
\]
第一行的反向蕴含在多变量、高次数时会失败，下面的 Motzkin 多项式给出反例。第二行是本节要证明的正性表示：允许有理函数的分母后，逐点非负性仍能转成平方和恒等式。

\begin{proposition}\label{b3:prop:univariate-sos}
设 \(R\) 为实闭域，\(f\in R[x]\)。则 \(f\) 在 \(R\) 上非负，当且仅当它是两个 \(R[x]\) 中多项式的平方之和。
\end{proposition}
\begin{proof}
平方和非负显然。对非零非负 \(f\)，实根的重数必须为偶数，否则根两侧的符号相反。实闭域上的因式分解把 \(f\) 写成正系数、偶次线性因子与正首项不可约二次因子之积。每个后一类因子配方为 \(c((x-a)^2+b^2)\)，其中 \(c>0,b\ne0\)，而正数在 \(R\) 中有平方根。再反复使用
\[
(u^2+v^2)(s^2+t^2)=(us-vt)^2+(ut+vs)^2
\]
即得结论；零多项式单独处理。
\end{proof}

第二卷定理 10.2.4 的 Sylvester 惯性定律还说明：任意变量数的实非负二次型都是线性形式的平方和。于是最先出现困难的不是“变量多”本身，而是多个变量与较高次数同时出现。

\subsection{Motzkin 多项式与 Artin 定理}
\label{b3:subsec:G03:02}

Motzkin 给出的多项式\footnote{莫茨金（Theodor Samuel Motzkin，1908-1970），德裔美国数学家，研究线性规划与凸几何，构造了非负但非多项式平方和的经典例子。}显示了多个变量与较高次数同时出现时的困难。

\begin{example}[Motzkin 多项式]\label{b3:exm:motzkin-not-sos}
多项式
\[
M(x,y)=x^4y^2+x^2y^4+1-3x^2y^2
\]
在 \(\mathbb R^2\) 上非负，却不是实多项式的平方和。
\end{example}
\begin{proof}
对三个非负数 \(x^4y^2,x^2y^4,1\) 用算术平均不小于几何平均，得其和不小于 \(3x^2y^2\)。若 \(M=\sum_jq_j^2\)，令 \(d\) 为各非零 \(q_j\) 的最大总次数。若 \(d>3\)，记 \(r_j\) 为 \(q_j\) 的 \(d\) 次齐次部分（次数较低者取 \(r_j=0\)），则右边的 \(2d\) 次齐次部分为 \(\sum_jr_j^2\)。左边的次数为六，故 \(\sum_jr_j^2=0\)。在每个实点求值，平方项非负迫使各 \(r_j\) 处处为零；实数域是无限域，处处为零的多项式必为零多项式，这与至少一个 \(r_j\ne0\) 矛盾。因此 \(d\le3\)。

置 \(y=0\) 得 \(\sum_jq_j(x,0)^2=1\)，同样比较最高次项便知每个 \(q_j(x,0)\) 为常数；置 \(x=0\) 同理。因此
\[
q_j=a_j+b_jxy+c_jx^2y+d_jxy^2.
\]
在 \(\sum_jq_j^2\) 中，\(x^2y^2\) 的系数只能是 \(\sum_jb_j^2\ge0\)，但 \(M\) 中该系数为 \(-3\)，矛盾。
\end{proof}

\begin{theorem}[Artin 定理]\label{b3:thm:artin-hilbert17}
设 \(n\ge1\)，\(f\in\mathbb R[x_1,\ldots,x_n]\)。若对每个 \(a\in\mathbb R^n\) 都有 \(f(a)\ge0\)，则 \(f\) 能在有理函数域 \(\mathbb R(x_1,\ldots,x_n)\) 中写成有限个平方之和。
\end{theorem}

证明将“不是平方和”转成某个序中的负性。反设 \(f\) 不是有理函数平方和，先证明由 \(-f\) 与全部平方生成的预序不含 \(-1\)；再把这个预序扩张为域序，使 \(f<0\)。第三步取保持该序的实闭包，于是 \(f<0\) 在一个实闭扩域中有解；第四步用量词消去把存在解的结论传回 \(\mathbb R\)，与原来的逐点非负性矛盾。

\begin{proof}
若 \(f=0\)，结论显然。设 \(f\ne0\)，记 \(L=\mathbb R(x_1,\ldots,x_n)\)、\(\Sigma L^2\) 为有限平方和。\(L\) 可以赋予延拓 \(\mathbb R\) 通常序的域序：逐次在 \(\mathbb R(x_1,\ldots,x_{i-1})(x_i)\) 中，按分子与分母关于 \(x_i\) 的最高次项系数的符号规定非零有理函数的符号。这与第一卷第 E.1.1 小节的有理函数域赋序方法相同。因此 \(-1\notin\Sigma L^2\)。若 \(f\notin\Sigma L^2\)，则由 \(-f\) 生成的预序
\(P=\Sigma L^2-f\Sigma L^2\) 不含 \(-1\)。事实上，若 \(-1=\sigma_0-f\sigma_1\)，其中 \(\sigma_0,\sigma_1\in\Sigma L^2\)，则 \(\sigma_1=0\) 会使 \(-1=\sigma_0\) 成为平方和，矛盾。故 \(\sigma_1\ne0\)，且
\[
f=\frac{1+\sigma_0}{\sigma_1}
=\frac{(1+\sigma_0)\sigma_1}{\sigma_1^2}
\]
是平方和，因为平方和的乘积仍为平方和，这与假设矛盾。

\(P\) 含有全部平方，对加法封闭；乘法封闭则由
\[
(\sigma_0-f\sigma_1)(\tau_0-f\tau_1)
=(\sigma_0\tau_0+f^2\sigma_1\tau_1)
-f(\sigma_0\tau_1+\sigma_1\tau_0)
\]
及平方和对和、积的封闭性得到。第一卷推论 E.1.4 的正锥扩张结论因而把 \(P\) 包含于 \(L\) 的一个域序的正锥。在这个序中 \(-f\) 为正；\(f\ne0\) 保证 \(f<0\)。再由第一卷定理 E.2.5，取保持该序的实闭包 \(L^{\mathrm{rc}}\)。其在 \(\mathbb R\) 上的序与通常序相同，因为正实数都是非零平方。于是带实系数的不等式 \(f(x)<0\) 在 \(L^{\mathrm{rc}}\) 中有解。\cref{b3:thm:real-projection}的量词消去把这个存在公式化为只涉及实系数的符号条件；这些条件在 \(L^{\mathrm{rc}}\) 与 \(\mathbb R\) 中相同，所以它在 \(\mathbb R\) 中也有解，违背 \(f\) 处处非负。因此 \(f\in\Sigma L^2\)。
\end{proof}

这是 Hilbert 第17问题的肯定解答，也可由正性证书推导，见\cite[Section 7.4, Corollary 7.6]{Sturmfels2002PolynomialEquations}。该结论允许分母，因而不与上例矛盾。

\subsection{平方和表示的系数域、分母与定义域}
\label{b3:subsec:G03:03}

把 Artin 定理中的各分母通分，可取非零多项式 \(q\in\mathbb R[x_1,\ldots,x_n]\) 与多项式 \(p_1,\ldots,p_s\)，使
\[
\boxed{q^2f=p_1^2+\cdots+p_s^2.}
\]
这是整个多项式环中的代数证书，可以在每个 \(a\in\mathbb R^n\) 代入。与它对应的分式表达只有在分母非零处才能逐项求值：
\[
f(a)=\sum_j\left(\frac{p_j(a)}{q(a)}\right)^2,
\qquad a\in\mathbb R^n,\quad q(a)\ne0.
\]
在 \(q(a)=0\) 处，多项式恒等式只直接给出各 \(p_j(a)=0\)，不能在这些点把分式代入，也不能仅由代入后的 \(0=0\) 判定 \(f(a)\) 的符号。对于实系数多项式，\(q\ne0\) 的集合在欧氏拓扑中稠密，故全局证书仍蕴含 \(f\) 处处非负：先在分母非零处判定符号，再用 \(f\) 的连续性延伸到分母零点。这里的稠密性来自非零多项式的零点集没有内点，可逐变量用一元多项式根数有限来证明。

\ProseParagraphBreak
系数域也应写清。\(2x^2\) 是一个实系数平方 \((\sqrt2x)^2\)，却不是一个有理系数多项式的平方；它仍可写成 \(x^2+x^2\)。这个例子说明“存在平方和”“存在规定数目的平方”“各平方的系数属于指定域”是不同问题。算法若以 \(\mathbb Q\) 为输入域，必须输出有理系数恒等式，或明确给出所用实代数数的最小多项式与隔离区间。浮点近似 \(1.4142\) 不能替代 \(\sqrt2\) 的精确身份。
